\section{Results and Analysis}
\label{sec:results}

\subsection{The metric is the first result}
\label{sec:res-metrics}

\paragraph{A zero-information baseline scores above the noise ceiling.} A predictor that places
every gene at the middle rank --- carrying no information about drug, cell, or biology --- scores
$\DEdr(K{=}50) = 0.9999$ [$0.99991$, $0.99992$], above the two-real-cell noise ceiling ($0.76$) and
above the model ($0.73$). The ordering is inverted: less information yields a higher score, which is
the signature of an artifact rather than of skill. Note also that this predictor's partial-$\DEdr$
and $\paneltau$ are \emph{undefined} rather than merely poor --- its shift from control is
identically zero, so there is nothing left to correlate once the control is regressed out. The cause is control regression-to-the-mean: $\DEdr$ correlates
predicted and true change from control, and a prediction that reverts toward the centre correlates
with the true change by construction.

\begin{table}[h]
\centering
\begin{tabular}{lrrr}
\toprule
Predictor & Information & $\DEdr$ (K50) & partial-$\DEdr$ \\
\midrule
\texttt{revert\_center} (all genes $\rightarrow P/2$) & none & \textbf{0.9999} & undefined \\
\texttt{revert\_mean} (predict mean control) & none, no fit & 0.961 & 0.25 \\
ridge linear (control $\rightarrow$ shift) & control & 0.947 & 0.28 \\
noise ceiling (two real cells) & --- & 0.76 & --- \\
model (1B LLM) & control + drug & 0.73 & $\approx 0.26$ \\
\bottomrule
\end{tabular}
\caption{$\DEdr$ is saturated by a control artifact. Regressing out the control leaves
$\approx0.26$ of real but \emph{drug-agnostic} skill, matched by a no-fit baseline.}
\label{tab:de-artifact}
\end{table}

Removing the control (partial-$\DEdr$) collapses the raw $0.95$ to $\approx 0.26$ --- and that
residue is matched by a baseline that performs no fitting at all. The response decomposes into
control reversion (an artifact), a generic drug program ($\approx 0.26$, drug-agnostic, trivially
matched), and a drug-specific component that no predictor captures.

\paragraph{Only one metric survives calibration.} Under \DRF{} computed within plate (61 groups,
655 drugs, $\approx53$ cells per drug):

\begin{table}[h]
\centering
\begin{tabular}{lrrrl}
\toprule
Metric & \DRF{} (within-plate) & $m$(neg) & $m$(pos) & Verdict \\
\midrule
$\NIR$          & \textbf{+0.446} & 0.274 & 0.597 & only calibrated metric \\
\texttt{weighted\_r2} & $-0.081$ & 0.793 & 0.776 & inverted \\
$\paneltau$     & $-0.231$ & 0.694 & 0.623 & inverted \\
\texttt{spearman\_expr} & $-0.415$ & 0.851 & 0.790 & inverted \\
$\DEdr$         & $-0.448$ & 0.886 & 0.834 & inverted \\
\bottomrule
\end{tabular}
\caption{A negative \DRF{} means a drug-agnostic baseline outscores the noise ceiling.}
\label{tab:drf}
\end{table}

The inversion is a property of the data, not a leak: a leave-one-out correction barely moves the
numbers, since real cell lines carry 20--204 drugs. It also survives holding the plate constant. The
one earlier claim that did not survive is retracted here: \texttt{weighted\_r2} is approximately
neutral rather than strongly inverted.

\paragraph{This replicates Miller et al.\ --- and then diverges from them.} Their central
methodological claim is confirmed emphatically: metric choice dominates, the rank-based metric is
the one that survives, and the correlation-based metrics they identify as poorly calibrated are
exactly the ones that fail here. The agreement extends to the specific metrics --- they find
$\NIR$ consistently calibrated across 14 Perturb-seq datasets, and we find it is the only
calibrated metric of the five we tested on drug perturbation. Our \texttt{weighted\_r2} result is
also consistent with their finding that weighted metrics behave better than unweighted ones.

The divergence comes at the next step. Miller et al.\ conclude that under calibrated metrics deep
models outperform uninformative baselines; Section~\ref{sec:res-blind} evaluates on the calibrated
metric and finds the model at chance. Two differences in setting are candidates for the
discrepancy --- genetic versus drug perturbation, and a discriminative regression model versus an
autoregressive generative one --- and Section~\ref{sec:res-lookup} supplies a third and more
structural possibility.

\paragraph{The metric can discriminate; the model cannot.} A spike-in forced choice between two real
drug populations separates them essentially perfectly at pseudobulk ($\paneltau = 1.000$, $\DEdr =
0.995$, top-$n$ $\tau = 0.987$, plate held constant). The instrument is not the bottleneck.

\subsection{The model is drug-blind}
\label{sec:res-blind}

\begin{table}[h]
\centering
\begin{tabular}{lrl}
\toprule
Instrument & Value & Reading \\
\midrule
Forced-choice grading & $0.48$ ($\approx$ scramble; ceiling 0.67--0.83) & chance \\
Output invariance (gap) & \ci{0.000}{-0.019}{+0.016} & null \\
Scramble $\DEdr$, tier 1 & $0.740$ real vs $0.739$ scrambled & identical \\
Scramble $\DEdr$, tier 2 & $0.733$ real vs $0.729$ scrambled & identical \\
$\gap$ (identifiable subset) & \ci{+0.014}{-0.016}{+0.042} & null \\
\bottomrule
\end{tabular}
\caption{Four independent instruments agree. Near-ceiling $\DEdr$ is achieved \emph{identically}
with a wrong-mechanism drug, so the score is drug-independent.}
\label{tab:blind}
\end{table}

On the calibrated metric the model sits at chance and is matched by predictors that know nothing
about the drug:

\begin{table}[h]
\centering
\begin{tabular}{lr}
\toprule
Predictor & $\NIR$ \\
\midrule
ceiling (real replicate) & 0.576 \\
model & 0.498 \\
linear (drug-agnostic) & 0.500 \\
control-copy (zero drug information) & 0.504 \\
mean & 0.180 \\
\bottomrule
\end{tabular}
\caption{Aggregate expression-space $\NIR$, within plate, tier 2 ($n=606$). Chance is $0.50$.}
\label{tab:nir-aggregate}
\end{table}

\paragraph{Two objections, both closed.} First, \emph{is the task winnable at all?} Only $46.2\%$ of
(drug, cell line) conditions are identifiable in principle; $27.3\%$ of drugs are statistically
inert against DMSO and $78.7\%$ have a plate-mate closer than their own replicate. But restricting
to the identifiable subset does not rescue the model: it appears to beat the linear baseline there
($0.768$ vs $0.531$), yet a zero-drug-information control-copy matches it exactly ($0.766$) and the
leak-immune $\gap$ remains null. The model is drug-blind \emph{even where the task is provably
winnable}. The ceiling itself rises with aggregation ($0.61$ at 10 cells, $0.76$ at 40, $0.85$ at
120), so the signal is real and noise-limited.

Second, \emph{does the scramble swap in similar drugs?} The swap-distance sweep resolves this:
$\gap$ is null in every distance bin including the far tail ($-0.019$), across a $2.3\times$ spread
in swap distance, with a trend correlation of $+0.051$. Telling the model it is a drug whose true
response is maximally different still does not change the output.

\subsection{Where the blindness lives}
\label{sec:res-mechanism}

Drug identity \emph{is} in the representation: linear probes recover it at $82\%$ (layer 9) and
$76\%$ (layer 16) against a chance rate of $\approx 8\%$. With cell line, plate and dose projected
out, decodability \emph{rises} with depth, from $0.354$ at layer 2 to $0.883$ at layer 16. Yet the
purified drug subspace's variance share stays flat at $\approx 0.03$, against $\approx 0.6$ for cell
line.

\begin{table}[h]
\centering
\begin{tabular}{lrrrrrr}
\toprule
Layer & Decode & Decode & Ablate pure & Random & Ratio & \textbf{Swap vs} \\
      & (raw)  & (context removed) & drug (KL) & matched dim & & \textbf{matched noise} \\
\midrule
2  & 0.408 & 0.354 & 0.155 & 0.012 & $12.7\times$ & \textbf{0.104 $<$ 0.917} \\
4  & 0.585 & 0.429 & 1.563 & 0.084 & $18.6\times$ & \textbf{1.855 $<$ 2.156} \\
6  & 0.692 & 0.529 & 1.435 & 0.251 & $5.7\times$  & \textbf{1.618 $<$ 1.958} \\
8  & 0.779 & 0.569 & 0.596 & 0.042 & $14.1\times$ & \textbf{0.599 $<$ 1.163} \\
9  & \textbf{0.821} & 0.602 & 0.418 & 0.100 & $4.2\times$  & \textbf{0.333 $<$ 0.761} \\
12 & 0.754 & 0.873 & 0.133 & 0.035 & $3.8\times$  & \textbf{0.105 $<$ 0.227} \\
16 & 0.758 & \textbf{0.883} & 0.015 & 0.004 & $3.6\times$  & \textbf{0.031 $<$ 0.036} \\
\bottomrule
\end{tabular}
\caption{All seven layers pass the removal gate (88--95\% of above-chance drug signal removed);
chance is $0.083$. Two decode series are shown because they behave differently and the distinction
matters: the \emph{raw} probe peaks at layer 9 and then declines, whereas with cell line, plate and
dose projected out decodability rises monotonically to layer 16. The swap is the identity test ---
drug-B injection against a random vector of identical norm, both with drug means estimated on a
held-out half --- and it is null at every depth. Swap values carry no dispersion: the probe reduces
the per-prompt distributions to two means before storing them.}
\label{tab:workspace}
\end{table}

\paragraph{The headline is the swap.} Overwriting drug A's code with drug B's, inside the purified
subspace, perturbs the output \emph{less} than matched-norm noise, at all seven depths. The readout
responds to the \emph{magnitude} of activity in the drug region, not to \emph{which} drug it
encodes. Held-out mean estimation reproduces this exactly (layer 2: $0.104$ vs $0.106$; layer 9:
$0.334$ vs $0.334$; layer 16: $0.031$ vs $0.031$), ruling out estimation noise as the cause.

\paragraph{Validity of the null.} Three pillars: the removal gate passes at every layer, so the
nulls are not broken hooks; the cell-line positive control fires at 15--20$\times$ the drug's
effect, so the instrument works; and the swap is translation-clean by matched norm.

\paragraph{Reading.} The clean claim is not that the drug lives in a dead subspace. It is that
\textbf{the generation readout is drug-agnostic within an increasingly explicit drug
representation}. This is a concrete instance of representational selectivity~\cite{gurnee}:
cell-line identity is read automatically because emitting any plausible profile requires it, while
drug$\rightarrow$gene response demands flexible, context-dependent composition and never enters the
causally active subspace. It also \emph{predicts} the next result --- if the defect is the readout
ignoring direction, no change to the target can touch it.

\subsection{The $\varepsilon$-ladder: no target-side fix works}
\label{sec:res-epsilon}

\begin{table}[h]
\centering
\begin{tabular}{llrl}
\toprule
Target & $\varepsilon$ & $\gap$ & Verdict \\
\midrule
single cell (real treated cell) & $\approx 0$ & \ci{+0.014}{-0.016}{+0.042} & null \\
OT barycentric (state-matched) & $0.5$ & \ci{-0.001}{-0.018}{+0.016} & null \\
consensus (global mean) & $\rightarrow\infty$ & \ci{-0.010}{-0.022}{+0.003} & null \\
\bottomrule
\end{tabular}
\caption{The regularisation parameter interpolates between the two target constructions tested
independently. Every point on the ladder is drug-blind.}
\label{tab:epsilon}
\end{table}

Training was not inert: the consensus model stops echoing the control ($0.653$ against a
control-copy of $0.766$, where the single-cell model sits \emph{on} the control at $0.771$) and its
predictions are not mode-collapsed. But the variation is uncorrelated with real drug identity
(Mantel $\approx +0.03$). Not collapse --- wrong structure.

As predicted by Section~\ref{sec:res-mechanism}, denoising changes \emph{what is in} the target
while the defect is the readout ignoring direction within it.

\subsection{The repair: it was the tokenisation}
\label{sec:res-encoding}

\begin{table}[h]
\centering
\begin{tabular}{lrr}
\toprule
Representation & Inter-drug top-200 tokens differing & Replicate retrieval ceiling \\
\midrule
full profile (the standard target) & \textbf{34.6 / 200 (17\%)} & 0.742 \\
shift (treated $-$ control) & 111.4 (56\%) & 0.742 \\
\textbf{residual (drug-specific)} & \textbf{120.3 / 200 (60\%)} & 0.747 \\
\bottomrule
\end{tabular}
\caption{$83\%$ of the standard target's tokens are identical between any two drugs in the same
context. Note the ceiling is \emph{unchanged}: the information was always present, and the
tokenisation discarded it.}
\label{tab:divergence}
\end{table}

Because a cell sentence ranks genes by expression, cross-entropy allocates roughly $1\%$ of its
gradient to drug identity. This explains why every change to target \emph{content} failed
identically. Re-encoding the target as the drug-specific residual raises the differing fraction to
$60\%$ and produces the first non-null result in the project:

\begin{table}[h]
\centering
\begin{tabular}{lrrr}
\toprule
Swap stratum & Residual model & Single-cell (control) & OT/T2 (control) \\
\midrule
\texttt{near} ($\cos = +0.41$) & \ci{+0.0351}{+0.0061}{+0.0656} & $-0.0140$ & $+0.0073$ \\
\texttt{orth} ($\cos \approx 0$) & \ci{+0.0716}{+0.0371}{+0.1059} & $-0.0206$ & $+0.0144$ \\
\textbf{\texttt{opposite}} ($\cos = -0.33$) & \ci{+0.1429}{+0.1112}{+0.1787} & $-0.0204$ & \ci{+0.0263}{+0.0069}{+0.0476} \\
\midrule
Gradient with dissimilarity & \textbf{yes (all three clear zero)} & no (flat) & yes \\
\bottomrule
\end{tabular}
\caption{Stratified scramble. The monotone gradient is itself the evidence: a metric artifact would
not order itself by swap dissimilarity. \todo{the two control columns ran at a 600-token generation
budget against the residual model's 1400 --- see the caveat below. Re-run before quoting the
cross-model comparison.}}
\label{tab:residual-strata}
\end{table}

\paragraph{A confound in the control columns, stated rather than buried.} The residual model's
figures come from a run at $1400$ generation tokens; both control arms ran at $600$. Since the token
budget alone doubles the measured effect for the residual model, the controls were given half the
treatment's budget. The bias runs in the direction that flatters our conclusion --- truncation drives
a gap toward zero, so an under-budgeted control looks \emph{more} null than it is --- which is
precisely why it must be stated. Mitigating but not sufficient: a cell sentence is roughly $123$ gene
symbols against a residual signature's $201$, so the cap binds far less on the controls; the
completion rates were logged but not persisted, so this cannot be verified from the artifacts.
\todo{re-run both controls at 1400} Two things are unaffected: the residual model's own gradient
across strata, every point of which clears zero within a single matched run, and all of
Section~\ref{sec:res-generalisation}, where every arm ran at $1400$ inside one job.

Three further checks agree, and only the residual model passes all three: the correlation between
condition reproducibility and model $\NIR$ is $+0.144$ (positive: the model does better where the
drug effect is real) against $-0.117$ and $-0.216$ for the controls; prediction diversity is highest
for the residual model; and $\cos$(prediction, own truth) exceeds $\cos$(prediction, other drugs)
only for it.

\paragraph{A methodological note.} The effect was first measured at a $600$-token generation cap,
which silently truncated $26\%$ of generations --- a signature is $\approx201$ gene symbols at 2--4
BPE tokens each --- and \emph{halved} it ($+0.0716$ against the correct $+0.1429$). For any
generation-based evaluation, log the completion rate.

\paragraph{A retraction.} The OT arm was originally reported as achieving nothing. Under the sharper
stratified scramble it scores \ci{+0.0263}{+0.0069}{+0.0476} with a monotone gradient: it does use
the drug, weakly. The original null was an underpowered instrument, not an absence of signal. The
$\varepsilon$-ladder conclusion stands, but ``OT achieved nothing'' is withdrawn.

\subsection{The drug use generalises}
\label{sec:res-generalisation}

\begin{table}[h]
\centering
\begin{tabular}{lrrrl}
\toprule
Split & $n$ (cell lines) & model $\NIR$ & $\gap$ (opposite swap) & Verdict \\
\midrule
\texttt{train} & 200 (40) & 0.657 & \ci{+0.0898}{+0.0530}{+0.1269} & uses drug \\
\textbf{\texttt{unseen\_combo}} & \textbf{250 (38)} & \textbf{0.650} & \ci{+0.1002}{+0.0661}{+0.1368} & \textbf{transfer} \\
\texttt{unseen\_drug} & 120 (36) & 0.586 & \ci{+0.0195}{-0.0452}{+0.0860} & null (by design) \\
\bottomrule
\end{tabular}
\caption{Three-way holdout keyed at the (drug, cell line) level. Validity clean first: 9,120
generations, \DOWN{} present in $93\%$, $98.9\%$ valid panel genes, $1.4\%$ duplicates.}
\label{tab:generalisation}
\end{table}

\paragraph{There is no memorisation premium.} Held-out pairings score \emph{as well as} trained ones
($+0.1002$ vs $+0.0898$, intervals almost entirely overlapping). The model did not memorise
pairings; it learned a drug representation that applies in cell lines where it never saw that drug.
The \texttt{unseen\_drug} arm behaving as designed --- same model, same instrument, drug token
carrying no learned information, no gap --- is what licenses attributing the transfer to having seen
the drug.

\paragraph{Two instrument defects, both of which would have concealed this.} A uniform condition
sampler reproduces the pool's split proportions, so $250$ constructed held-out combinations yielded
only $33$ scored ones; and a strict-inequality $\NIR$ scored the zero vector (``predict the generic
response'') at $0.000$ rather than $0.500$. Both are corrected in Section~\ref{sec:methods}, and
after correction the drug-agnostic baselines land where they must: \texttt{generic} exactly
$0.500$, \texttt{control\_copy} $0.478$, \texttt{random} $0.502$. The residual frame is therefore
\emph{demonstrably} leak-proof rather than assumed to be.

\subsection{But a lookup table wins}
\label{sec:res-lookup}

\begin{table}[h]
\centering
\begin{tabular}{lrl}
\toprule
Arm & $\NIR$ & model $-$ arm (clustered CI) \\
\midrule
ceiling (within-condition replicate) & \textbf{0.968} & --- \\
\texttt{drug\_lookup} & \textbf{0.963} & \ci{-0.3245}{-0.3539}{-0.2975} loses \\
\texttt{drug\_lookup\_1} & \textbf{0.832} & \ci{-0.1936}{-0.2285}{-0.1612} loses \\
model & 0.639 & --- \\
\texttt{moa\_lookup} ($n=143$) & 0.529 & \ci{+0.1042}{+0.0144}{+0.1882} wins \\
\texttt{generic} & 0.500 & \ci{+0.1388}{+0.1086}{+0.1650} wins \\
\texttt{random} & 0.502 & --- \\
\texttt{control\_copy} & 0.478 & \ci{+0.1613}{+0.1314}{+0.1907} wins \\
\bottomrule
\end{tabular}
\caption{The baseline ladder, $n=570$. On the held-out split specifically, model $-$
\texttt{drug\_lookup} is $-0.3170$ and model $-$ \texttt{drug\_lookup\_1} is $-0.1825$.}
\label{tab:lookup}
\end{table}

The pre-registered interpretation, written before the result, anticipated three branches; the middle
one fired. The model reads and transfers the drug, but recovers only a small fraction of its
signature. That it beats \texttt{moa\_lookup} rules out the explanation that it is simply reading the
\texttt{Mechanism:} field of the prompt --- though at $0.529$ that is a low bar.

\paragraph{This is DrEval's finding, in transcriptomic space.} DrEval reports that deep learning
models barely outperform a naive predictor of mean drug and cell-line
effects~\cite{dreval}. Their \texttt{NaiveMeanEffectsPredictor} is
$\hat{y}_{ij} = \mu^c_i + \mu^d_j - \mu$; our \texttt{drug\_lookup} is the $\beta(d)$ term of that
same additive model, measured in a frame where the $\textrm{generic}(c)$ term has already been
subtracted. The finding replicates and sharpens: on a $946$-gene transcriptomic response rather than
a scalar potency, and with the comparison made against a proper within-condition replicate ceiling,
the additive drug main effect reaches $0.963$ of a $0.968$ ceiling while a $1$B-parameter model
reaches $0.639$.

Note also that this result lands in the split where DrEval specifically warns the drug-average
baseline will look strong: \texttt{unseen\_combo} is Leave-Pairs-Out, and DrEval reports
that good LPO performance is reachable by predicting each drug's training average~\cite{partin}.
We take that warning as the hypothesis rather than the caveat, and measure it. It holds.

\paragraph{A retraction, made before it entered the record.} It is tempting to read
$\textrm{ceiling} - \texttt{drug\_lookup} = +0.005$ as ``the modelling target has vanished''. That
inference is not supported, for three reasons that all favour the lookup: the ceiling is scored
against a \emph{half-sample} truth while every baseline is scored against the \emph{full-sample}
truth; \texttt{drug\_lookup} averages $\approx38$ conditions against the ceiling's single noisy half;
and $\NIR$ near $0.96$ is compressive, collapsing a wide range of cosines into a few thousandths.

\subsection{How much is left? The transfer coefficient}
\label{sec:res-transfer}

Sections~\ref{sec:res-generalisation} and~\ref{sec:res-lookup} leave two numbers pointing opposite
ways, and $\NIR$ cannot adjudicate between them. \texttt{drug\_lookup\_1} sits $0.136$ below the
ceiling despite using \emph{more} cells than the ceiling's half-sample, arguing for real
cell-line-specific structure; yet averaging many lines recovers to $0.963$, and averaging removes only
\emph{noise}, never the systematic miss of the target line's own interaction. The question is
therefore posed in cosine space with the noise divided out (Section~\ref{sec:methods-vardecomp}).

\paragraph{A false alarm, and the control-design lesson it produced.} The first run compared
$\Tcoef$ against a negative control of \emph{different drugs on the same plate}, which returned
$+0.485$ against $\Tcoef = 0.511$, and was correctly declared void. But that control differs from
$\Tcoef$ in \textbf{two} ways at once --- different drug \emph{and} same cell line. Plate structure
surviving a cell-line-scoped generic inflates \emph{same-plate} comparisons while leaving
\emph{cross-line} comparisons untouched, so it destroyed the control's usefulness without ever
touching $\Tcoef$. A \emph{structure-matched} control --- different drug, \textbf{different cell
line}, breaking only the drug match --- resolves it, and reads $+0.000$ in every configuration. The
episode is worth reporting rather than hiding: a negative control that differs from the estimand in
more than one way cannot bound it.

\begin{table}[h]
\centering
\begin{tabular}{lrrrrr}
\toprule
Configuration & $\Tcoef$ & same-plate null & \textbf{matched null} & dose & shared$-$split \\
\midrule
\textbf{plate-scoped generic} & \textbf{0.557} [0.513, 0.601] & $-0.018$ & \textbf{+0.000} & 0.703 & $-0.000$ \\
plate + generic projected out & 0.545 [0.503, 0.585] & $-0.019$ & +0.010 & 0.682 & $-0.004$ \\
cell-line-scoped (production) & 0.517 [0.474, 0.560] & $+0.478$ & +0.000 & 0.484 & $+0.254$ \\
\bottomrule
\end{tabular}
\caption{The transfer coefficient. Against a simulated $\kappa = 0$ null of $1.001$, the shortfall is
$0.444$--$0.484$. The matched negative control is $\approx 0$ throughout.}
\label{tab:transfer}
\end{table}

\paragraph{Answer: roughly 45\% of the drug-specific residual variance is drug$\times$cell-line
interaction}, robust to the generic's scope, to projecting the generic direction out, and to control
splitting. The per-drug constant $\beta(d)$ --- which \texttt{drug\_lookup} reproduces at $0.963$ ---
is therefore \emph{not} the whole signal. There is a conditional target worth about half the
drug-specific variance that no lookup can structurally reach, and that the model does not reach
either.

Three internal checks agree. \textbf{Dose ordering}: at plate scope, same-drug/same-line/different-dose
transfer is $0.703$, \emph{above} cross-cell-line $0.557$ --- changing the dose costs less than
changing the cell line, the biologically correct ordering; at cell-line scope the ordering inverts.
\textbf{Control-split bias}: at plate scope the shared- and split-control builds agree to $-0.000$,
against $+0.254$ at cell-line scope, so that discrepancy was plate contamination rather than the split.
\textbf{Orthogonality}: $(\lVert\resid\rVert/\lVert\shift\rVert)^2 +
(\lVert\textrm{generic}\rVert/\lVert\shift\rVert)^2 = 1.002$, so the two components are essentially
orthogonal and the fractions below are genuine variance shares.

\paragraph{Scope of the whole thesis.} At plate scope $\lVert\resid\rVert/\lVert\shift\rVert = 0.786$
and $\lVert\textrm{generic}\rVert/\lVert\shift\rVert = 0.620$: \textbf{62\% of the perturbation
response variance is drug-specific}, and 38\% is the generic program plus batch. Every claim in this
thesis is scoped by this number, and it has not previously been computed for this task.

\paragraph{A retraction this forces.} The justification for estimating the generic per cell line
rather than per plate was a reproducibility comparison, $62\%$ against $19\%$. That was measured with
\emph{shared} control cells, which inflate cell-line reproducibility far more than plate
reproducibility. Measured with split controls the two are comparable ($20\%$ plate, $16\%$ cell line),
so the argument does not favour cell-line scope. Combined with the dose-ordering inversion and the
control-split bias, the plate-scoped build is the better one --- and the residual training targets,
built at cell-line scope, therefore retain batch structure. This does not affect $\gap$, which is a
paired contrast in which any shared component cancels, but it makes a plate-scoped rebuild the
natural next arm.

\subsection{Scope: what is closed by measurement}
\label{sec:res-scope}

\begin{table}[h]
\centering
\begin{tabular}{llr l}
\toprule
Channel & Measurement & Value & Verdict \\
\midrule
chemical structure $\rightarrow$ response & Morgan / MolFormer retrieval & 0.500 / 0.474 & closed (ceiling 0.805) \\
mechanism $\rightarrow$ response & \texttt{moa\_lookup} & 0.529 & closed \\
mechanism $\rightarrow$ response & within/between MoA distance & ratio 0.977 & closed \\
\textbf{protein target $\rightarrow$ response} & --- & --- & \textbf{untested} \\
\bottomrule
\end{tabular}
\caption{An unseen drug can only be predicted through a property it shares with drugs seen in
training. Two such channels are closed by measurement; one has never been tested.}
\label{tab:channels}
\end{table}

Molecular structure does not predict response here, which \emph{scopes} the whole project: the
positive results of Sections~\ref{sec:res-encoding} and~\ref{sec:res-generalisation} are seen-drug
results. Note carefully what is \emph{not} established: the structure gate tested chemical
fingerprints only, and mechanism was tested at the coarse \texttt{moa-fine} granularity. Protein
target --- a finer relation, and the one a knowledge-based approach would exploit --- remains
untested. \todo{\texttt{channel\_gate.py}}
